\chapter{Nephrectomy}
\section*{Overview}
Nephrectomy removes part of a kidney or the entire kidney. Partial nephrectomy preserves functioning renal tissue when feasible.
\section*{Why This Test or Procedure Is Ordered}
It may be performed for kidney cancer, severe trauma, a nonfunctioning infected kidney, donation, or selected complex renal disease.
\section*{How This Test or Procedure Is Performed}
The operation may be open, laparoscopic, or robotic. The renal vessels and ureter are controlled, and the kidney or diseased portion is removed.
\section*{Risks and Recovery}
Bleeding, infection, injury to nearby organs, urine leak, blood clots, reduced kidney function, and need for dialysis are possible. Kidney function and blood pressure require follow-up.
\section*{References}
\begin{enumerate}\item MedlinePlus. Kidney Removal. U.S. National Library of Medicine; 2025.\item Current Medical Diagnosis and Treatment. McGraw-Hill; 2025.\end{enumerate}
\clearpage
